\scanpage{135}
\mspage{80}
% archive sheet number: 134

C'est ce qui est obtenu à l'aide du foncteur ``montant'' (du \unsure{r. de ch.} vers l'étage) $C$, qu'on peut regarder comme un foncteur $\mathrm{Cat}_{/S}\unsure{\mathrel{\relbar\joinrel\relbar\joinrel\vert}|\longrightarrow}\underline{\mathrm{Fibhot}}\,S$, mais qui en fait prend ses valeurs dans $\underline{\mathrm{Parf}}\,S$. (C'est là un détail sans grande importance, me semble-t-il.)

\[
\mathrm{Cat}_{/S}\xrightarrow{C}\underline{\mathrm{Parf}}\,S\subset\underline{\mathrm{Fibhot}}\,S.
\]

Ce foncteur \deleted{\unsure{transforme}} est compatible avec les localiseurs $W_S$ des deux côtés, \deleted{\unsure{donc}} \inserted{et} établit des équivalences (quasi-inverses des équivalences

\[
(\underline{\mathrm{Parf}}\,S)\deleted{\unsure{\Sigma}}W_S^{-1}\xrightarrow{\approx}(\mathrm{Cat}_{/S})W_S^{-1}
\quad\text{et}\quad
(\underline{\mathrm{Fibhot}}\,S)W_S^{-1}\overset{\approx}{\unsure{\mathrel{\relbar\joinrel\relbar\joinrel\vert}|\longrightarrow}}(\mathrm{Cat}_{/S})W_S^{-1}\bigr)
\]

\[
(\mathrm{Cat}_{/S})W_S^{-1}\xrightarrow{\approx}(\underline{\mathrm{Parf}}\,S)W_S^{-1}
\bigl(\overset{\approx}{\unsure{\mathrel{\relbar\joinrel\relbar\joinrel\vert}|\longrightarrow}}(\underline{\mathrm{Fibhot}}\,S)W_S^{-1}\bigr).
\]

Donc, quand les trois catégories $\mathrm{HOT}$ pour $S$ sont interprétées en termes de localisées de $\mathrm{Cat}_{/S}$, c'est le foncteur $C$ \inserted{(plus \unsure{exactemt|exactement} le composé $\mathrm{Cat}_{/S}\xrightarrow{C}\underline{\mathrm{Parf}}\hookrightarrow\mathrm{Cat}_{/S}$)} qui fournit à la fois le foncteur d'inclusion (!) $i$, et le foncteur d'inclusion $j$ dans $\mathrm{HOT}_{W}(S)$ et $\mathrm{HOT}_{W}^{\unsure{\circ|o}}(S)$ respectivement. Le foncteur en question

\[
C':\mathrm{Cat}_{/S}\longrightarrow\mathrm{Cat}_{/S}
\]

satisfait bien

\[
C'(W_S)\subset W_S^d\cap W_S^g
\]

et définit les deux foncteurs

\diagram{\begin{tikzcd}[ampersand replacement=\&,row sep=0.7em,column sep=large]
\&(\mathrm{Cat}_{/S})(W_S^d)^{-1}\approx\mathrm{HOT}_W(S)\\
{\begin{gathered}(\mathrm{Cat}_{/S})W_S^{-1}\\\approx\mathrm{HOT}_{W}^{\unsure{\mathrm{lc}|\mathrm{le}}}(S)\end{gathered}}\arrow[ur,"i"]\arrow[dr,dash,"j"']\&\\
\&(\mathrm{Cat}_{/\unsure{J|S}})(W_S^g)^{-1}\approx\mathrm{HOT}_{W}^{\circ}(S),
\end{tikzcd}}

Cela nous permet d'expliciter les formules d'adjonction de \textcircled{IV}, quand on y interprète $\unsure{\xi|x},\unsure{y|\eta}$
